\DocumentMetadata{
  pdfversion=2.0,
  pdfstandard=ua-2,
  lang=en-US,
  tagging=on,
  tagging-setup={math/setup=mathml-SE,tabsorder=structure}
}
\documentclass[11pt,letterpaper]{article}
\usepackage[margin=0.68in,includefoot]{geometry}
\usepackage{newcomputermodern}
\usepackage{amsmath,amscd,mathtools}
\usepackage{tikz,adjustbox}
\providecommand{\square}{\mdlgwhtsquare}
\usepackage[hidelinks]{hyperref}
\hypersetup{
  pdftitle={Algebra First-Year Exam, January 2024, Part 1},
  pdfauthor={University of Florida Department of Mathematics},
  pdfsubject={Graduate mathematics examination},
  pdfdisplaydoctitle=true,
  bookmarksopen=true
}
\setcounter{secnumdepth}{0}
\setlength{\parindent}{0pt}
\setlength{\parskip}{0.28em}
\setlength{\emergencystretch}{3em}
\setlength{\leftmargini}{2.15em}
\setlength{\leftmarginii}{2.65em}
\setlength{\leftmarginiii}{3em}
\renewcommand{\labelenumi}{\textbf{\theenumi.}}
\pagestyle{empty}
\raggedbottom
\allowdisplaybreaks
\newenvironment{examproblems}[1]{%
  \begin{enumerate}%
  \setcounter{enumi}{#1}%
  \setlength{\topsep}{0.35em}%
  \setlength{\partopsep}{0pt}%
  \setlength{\itemsep}{0.48em}%
  \setlength{\parsep}{0.18em}%
}{\end{enumerate}}
\newenvironment{examsubparts}{%
  \begin{enumerate}%
  \setlength{\topsep}{0.18em}%
  \setlength{\partopsep}{0pt}%
  \setlength{\itemsep}{0.16em}%
  \setlength{\parsep}{0.1em}%
}{\end{enumerate}}
\newenvironment{examinstructions}{%
  \par\begingroup\itshape
}{%
  \par\endgroup\vspace{0.18em}
}
\makeatletter
\renewcommand\section{\@startsection{section}{1}{0pt}%
  {-1.25ex plus -.25ex minus -.1ex}%
  {0.55ex plus .1ex}%
  {\normalfont\large\bfseries}}
\renewcommand{\@maketitle}{%
  \newpage\null\vskip -1em%
  {\footnotesize\bfseries
    \href{https://www.ufl.edu/}{University of Florida}, \href{https://math.ufl.edu/}{Department of Mathematics}\par}%
  \vskip -0.5em%
  \tagstructbegin{tag=Title}%
    {\LARGE\bfseries\href{https://gma.math.ufl.edu/exams/algebra-first-year/}{Algebra First-Year Exam}, January 2024, Part 1\par}%
  \tagstructend%
  \vskip 0.25em%
  \tagmcbegin{artifact}%
    \hrule height 1pt%
  \tagmcend%
  \vskip 1em%
}
\makeatother
\title{Algebra First-Year Exam, January 2024, Part 1}
\author{}
\date{}
\begin{document}
\maketitle
\thispagestyle{empty}
\begin{examinstructions}
Answer four problems. (If you turn in more, the first four will be graded.) Within reason, you may quote theorems as long as you state them clearly.
\end{examinstructions}
\begin{examproblems}{0}
\item
State and prove Sylow’s First Theorem.
\item
Let~\(G\) and~\(H\) be finite groups of orders~\(m\), and~\(n\) respectively. Justify your claims.
\begin{examsubparts}
\item[\textnormal{(a)}]
Prove that if~\(m\) and~\(n\) are coprime then
\[
\operatorname{Aut}(G \times H) \simeq \operatorname{Aut}(G) \times \operatorname{Aut}(H).
\]
\item[\textnormal{(b)}]
Give an example where~\(m\) and~\(n\) are not relatively prime and
\[
\operatorname{Aut}(G \times H) \not\simeq \operatorname{Aut}(G) \times \operatorname{Aut}(H).
\]
\end{examsubparts}
\item
Here~\(\mathbb{Z}\) is considered as a group. Show that there are two different group homomorphisms \(\phi : \mathbb{Z}/2\mathbb{Z} \to \operatorname{Aut}(\mathbb{Z}/7\mathbb{Z})\), and prove that the corresponding semidirect products \(\mathbb{Z}/7\mathbb{Z} \rtimes_{\phi} \mathbb{Z}/2\mathbb{Z}\) are pairwise nonisomorphic.
\item
Note that \(2024 = 2^3 * 11 * 23\).
\begin{examsubparts}
\item[\textnormal{(a)}]
Prove that the symmetric group \(S_{42}\) has at least one element of order \(2024\).
\item[\textnormal{(b)}]
Does the alternating group \(A_{42}\) have an element of order \(2024\)? Justify your answer.
\end{examsubparts}
\item
Let~\(G\) be a finite group.
\begin{examsubparts}
\item[\textnormal{(a)}]
Define what it means to say that~\(G\) is nilpotent.
\item[\textnormal{(b)}]
Without assuming any result about nilpotent groups (i.e. just from your definition), prove that if the order of~\(G\) is \(p^n\), where~\(p\) is a prime and \(n \in \mathbb{Z}\) with \(n \geq 0\), then~\(G\) is nilpotent.
\end{examsubparts}
\item
Let~\(S\) be a finite simple group, and let~\(H\) be a proper subgroup of~\(S\). Let \(n = [S : H]\). Prove that~\(S\) is isomorphic to a subgroup of the symmetric group \(S_n\).
\end{examproblems}
\end{document}
